\documentclass[12pt]{article}

\newcommand{\gname}{Landn{\aa}m}
\title{\gname Concepts}

\begin{document}
\maketitle
\section{Overview}

\gname is a token-placing game, where the tokens represent either
units of population (POPs) or status effects on them. Players compete
to have their POPs gain score in four categories: Consumption,
investment, defense, and meaning.

\section{Map}

The game takes place on a triangle-tile map; each tile has ten slots,
which may contain POPs, buildings, capital equipment, and
infrastructure. POPs may move on the map in particular circumstances
but most POPs will remain in one place for most of the game.

\section{POPs}

The population units represent a number of households, perhaps as many
as 100 people on average. They come in classes such as Peasant,
Knight, and Merchant, each of which produces different goods and
status effects and interacts with other POPs in different ways.

POPs contain ``surplus population'' and grow in size every turn; if
the surplus grows sufficiently, another POP is produced. 

\section{Production}

The base unit of production is the 'prod'. Without upgrades and
modifiers, a basic Peasant produces enough prods to support itself and
one-eighth of another POP. Each upgrade level of the Peasant doubles
the surplus (not total!) production. Prods can be transferred both
inside and outside of the producing POP's triangle, the possible
distance depending on what infrastructure (e.g. roads, harbours) and
what transferring POPs (Carters, Merchants) are present.

\section{POP ideas}

\begin{itemize}
\item Peasant (upgrade to Tenant, Smallholder, Yeoman). Basic
  production unit.
\item Bandit. Attempts to steal prods to survive.
\item Knight (upgrade to Lord, Duke, King). Defender: Increases
  militia efficiency again Bandits, basic military unit.
\item Smith. Increases capital investment.
\item Preacher (upgrade to Priest, Bishop, Archbishop). Increases
  meaning; needed for large-scale status effects.
\item Mason - creates buildings and roads.
\item Shipwright - allows water transport.
\item Carter - allows short-distance trade.
\item Merchant - allows trade over great distance.
\end{itemize}

\section{Status effect ideas}

\begin{itemize}
\item Be Fruitful and Multiply - increase birth rate.
\item Work Hard, Pray Hard - increase preacher effect, production. 
\item Drill Baby Drill - increase militia effect.
\item Best Time To Plant A Tree - increase capital investment.
\item Eat Drink And Be Merry - increase consumption.
\item Monopoly - excludes all but one Merchant, who can then extract
  more from the trade.
\end{itemize}

\section{Trade}

POPs can specialise; this is done by placing a merchant, who then
causes his whole triangle to specialise in a chosen good. POPs with
different specialisations can increase their production by trading. The
merchant goes looking for places to trade with; it must have a
different specialisation. If he finds one, both parties gain a
specialisation bonus to what they produced, which increases with every
trade, and also a specialisation penalty to the thing traded for,
which likewise increases. The bonus is larger than the penalty, such
that trade increases net production and repeated trade even more
so. But, if the trade does not go through, the POP still takes the
penalty but does not gain its bonus - reducing production possibly
drastically when trade is disrupted! Trade with unspecialised POPs
also possible but not as advantageous.

Advantages:

Merchant placement is now a crucial skill
Bandits, armies, and toll roads can all disrupt trading networks and must be managed
Automatic economic growth just from keeping the peace, but highly disruptible without additional tracking
Patterns of sustainable specialisation and trade, yay
Achievements: Long-distance trade, undisrupted trade ("Virgin With a Bag of Gold" cheevo)
Infrastructure gives additional options organically - roads, harbours - can now be maintained by the traders or by the state, different strategies
Pirates and storms
Monopolies and privileges to allow the merchant to extract more of the surplus.
Likewise implement basic supply-and-demand by skewing the extra
production towards the rarer good.

\section{Turn order}

\begin{itemize}
\item Player action - order a POP or place a tile.
\item Production.
\item Raids and combat.
\item Trade.
\item Investment.
\item Consumption.
\item Growth and upgrades.
\end{itemize}

\section{Capital}

Prods may be converted into capital. POPs with a sufficient amount of
capital can upgrade. Capital can be lost in combat or by being
consumed in emergencies.

\section{Buildings}

Roads, walls, harbours.

\section{Game opening}

The player deploys between six and thirty POPs, concentrated or spread
out; they start as Peasants.
